% !TeX root = ../main.tex

\chapter{Rechnen mit ganzen Zahlen}

Bisher haben wir die Grundrechenarten als elementare
Rechenoperationen betrachtet, die in konstanter Zeit
ausgeführt werden können.

Für beliebig große Zahlen ist diese Annahme jedoch
nicht realistisch. Wir untersuchen daher die Laufzeit
arithmetischer Operationen in Abhängigkeit von der
Anzahl der Binärstellen.

Als \textbf{elementare Rechenschritte} betrachten wir
Operationen und Vergleiche auf Zahlen beschränkter Größe.


\section{Addition und Subtraktion}

Die Addition zweier Zahlen mit höchstens $l$ Binärstellen
kann mit der schriftlichen Additionsmethode durchgeführt
werden.

Dabei wird jede Stelle höchstens eine konstante Anzahl
von Malen betrachtet. Somit beträgt die Laufzeit $O(l)$.

Durch die Verwendung der Vorzeichen- oder
Komplementdarstellung lässt sich dieses Verfahren
auch auf negative Zahlen und die Subtraktion erweitern.

\thm{Laufzeit von Addition und Subtraktion}{
  Für zwei ganze Zahlen $x,y\in\bbZ$ können
  die Summe $x+y$ und die Differenz $x-y$
  in Laufzeit
  \[
    \boxed{O(l)}
  \]
  berechnet werden, wobei
  \[
    l:=1+\left\lfloor
      \log_2\bigl(\max\{|x|,|y|,1\}\bigr)
    \right\rfloor.
  \]
}

\begin{myproof}
  Bei der schriftlichen Addition werden die
  Binärstellen von rechts nach links durchlaufen.

  Pro Stelle sind nur konstant viele Operationen
  erforderlich. Daher beträgt die Laufzeit $O(l)$.

  Die Subtraktion kann durch Addition der
  negierten Zahl durchgeführt werden:
  \[
    x-y=x+(-y).
  \]

  Die Negation mittels Zweierkomplement
  benötigt ebenfalls höchstens $O(l)$ Schritte.

  Somit können beide Operationen in
  $O(l)$ Zeit durchgeführt werden.
\end{myproof}


\section{Multiplikation}

Die gewöhnliche schriftliche Multiplikation
zweier Zahlen mit höchstens $l$ Stellen
benötigt $O(l^2)$ Rechenschritte.

Mit dem \textbf{Algorithmus von Karatsuba}
lässt sich diese Laufzeit verbessern.

\subsection{Karatsubas Algorithmus}

Wir zerlegen zwei natürliche Zahlen $x,y$
in jeweils zwei kleinere Zahlen:
\[
  x=x'B+x'',
  \qquad
  y=y'B+y'',
\]
wobei $B$ eine Potenz der verwendeten Basis ist.

Dann gilt
\begin{align*}
  xy
  &=(x'B+x'')(y'B+y'')\\
  &=x'y'B^2+(x'y''+x''y')B+x''y''.
\end{align*}

Eine direkte Berechnung würde vier
Multiplikationen kleinerer Zahlen benötigen.

Karatsubas Idee besteht darin, nur
\textbf{drei Multiplikationen} auszuführen.

Wir definieren
\begin{align*}
  p&:=x'y',\\
  q&:=x''y'',\\
  r&:=(x'+x'')(y'+y'').
\end{align*}

Wegen
\[
  r=p+x'y''+x''y'+q
\]
folgt
\[
  x'y''+x''y'=r-p-q.
\]

Somit erhalten wir
\[
  \boxed{
    xy=pB^2+(r-p-q)B+q.
  }
\]

\nt{
  Der entscheidende Vorteil von Karatsubas
  Algorithmus liegt darin, dass nur drei
  statt vier rekursive Multiplikationen
  benötigt werden.

  Die zusätzlichen Additionen und
  Subtraktionen sind vergleichsweise günstig.
}


\dfn{Karatsubas Algorithmus}{
  \textbf{Eingabe:} $x,y\in\bbN$ in Binärdarstellung.

  \textbf{Ausgabe:} Das Produkt $xy$.

  \begin{enumerate}
    \item Falls $x<8$ und $y<8$,
      berechne $xy$ direkt.

    \item Andernfalls setze
      \[
        l:=1+\lfloor\log_2(\max\{x,y\})\rfloor,
        \qquad
        k:=\left\lfloor\frac l2\right\rfloor,
        \qquad
        B:=2^k.
      \]

    \item Zerlege die Zahlen:
      \begin{align*}
        x'&:=\left\lfloor\frac xB\right\rfloor,
        &x''&:=x\bmod B,\\
        y'&:=\left\lfloor\frac yB\right\rfloor,
        &y''&:=y\bmod B.
      \end{align*}

    \item Berechne rekursiv
      \begin{align*}
        p&:=x'y',\\
        q&:=x''y'',\\
        r&:=(x'+x'')(y'+y'').
      \end{align*}

    \item Gib das Ergebnis
      \[
        xy=pB^2+(r-p-q)B+q
      \]
      aus.
  \end{enumerate}
}


\ex{Karatsubas Algorithmus}{
  Wir berechnen $13\cdot11$ mit $B=4$.

  Zunächst zerlegen wir:
  \begin{align*}
    13&=3\cdot4+1,\\
    11&=2\cdot4+3.
  \end{align*}

  Also gilt
  \[
    x'=3,\quad x''=1,\qquad
    y'=2,\quad y''=3.
  \]

  Wir berechnen:
  \begin{align*}
    p&=3\cdot2=6,\\
    q&=1\cdot3=3,\\
    r&=(3+1)(2+3)=20.
  \end{align*}

  Damit folgt
  \begin{align*}
    13\cdot11
    &=6\cdot4^2+(20-6-3)\cdot4+3\\
    &=96+44+3\\
    &=143.
  \end{align*}
}


\thm{Laufzeit von Karatsubas Algorithmus}{
  Die Multiplikation zweier natürlicher Zahlen
  $x,y$ in Binärdarstellung ist mit
  Karatsubas Algorithmus in Laufzeit
  \[
    \boxed{O(l^{\log_2 3})}
  \]
  möglich, wobei
  \[
    l:=1+\lfloor\log_2(\max\{x,y\})\rfloor.
  \]
}

\begin{myproof}
  \textbf{Korrektheit:}

  Die Korrektheit folgt unmittelbar aus
  \[
    xy=pB^2+(r-p-q)B+q.
  \]

  \medskip
  \textbf{Laufzeit:}

  Sei $T(l)$ die maximale Anzahl elementarer
  Rechenschritte für Eingaben mit höchstens
  $l$ Binärstellen.

  Die drei rekursiven Multiplikationen
  erhalten Eingaben mit höchstens
  \[
    \left\lceil\frac l2\right\rceil+1
  \]
  Binärstellen.

  Alle übrigen Operationen benötigen
  insgesamt $O(l)$ Zeit.

  Somit existiert eine Konstante $c>0$ mit
  \[
    T(l)\le
    \begin{cases}
      c,&l\le3,\\[2mm]
      cl+3T\left(
        \left\lceil\frac l2\right\rceil+1
      \right),&l\ge4.
    \end{cases}
  \]

  Wir zeigen zunächst durch Induktion über
  $m\in\bbN_0$, dass
  \[
    T(2^m+2)
    \le c(4\cdot3^m-2\cdot2^m-1).
  \]

  \textbf{Induktionsanfang:}

  Für $m=0$ gilt
  \[
    T(3)\le c
    =c(4-2-1).
  \]

  \textbf{Induktionsschritt:}

  Angenommen, die Aussage gilt für $m-1$.

  Dann erhalten wir
  \begin{align*}
    T(2^m+2)
    &\le c(2^m+2)
      +3T(2^{m-1}+2)\\
    &\le c(2^m+2)
      +3c(4\cdot3^{m-1}
      -2\cdot2^{m-1}-1)\\
    &=c(4\cdot3^m-2\cdot2^m-1).
  \end{align*}

  Damit ist die Abschätzung bewiesen.

  Für beliebiges $l>2$ wählen wir
  \[
    m:=\lceil\log_2(l-2)\rceil.
  \]

  Wegen $l\le2^m+2$ folgt, unter Verwendung
  der monotonen Worst-Case-Schranke,
  \begin{align*}
    T(l)
    &\le T(2^m+2)\\
    &<4c\cdot3^m\\
    &<12c\cdot3^{\log_2 l}\\
    &=12c\cdot l^{\log_2 3}.
  \end{align*}

  Also gilt
  \[
    T(l)\in O(l^{\log_2 3}).
  \]
\end{myproof}

\nt{
  Da
  \[
    \log_2 3\approx1{,}585<2,
  \]
  ist Karatsubas Algorithmus asymptotisch
  schneller als die schriftliche Multiplikation:
  \[
    O(l^{1{,}585})
    \quad\text{statt}\quad
    O(l^2).
  \]

  In der Praxis werden kleine Zahlen direkt
  multipliziert. Erst bei größeren Zahlen
  lohnt sich die rekursive Zerlegung.
}


\section{Euklidischer Algorithmus}

Der \textbf{euklidische Algorithmus}
berechnet den größten gemeinsamen Teiler
zweier natürlicher Zahlen.

Insbesondere kann man damit Brüche kürzen.


\dfn{Größter gemeinsamer Teiler und kgV}{
  Für $a,b\in\bbN$ bezeichnet
  \[
    \operatorname{ggT}(a,b)
  \]
  den größten gemeinsamen Teiler von $a$ und $b$.

  Falls
  \[
    \operatorname{ggT}(a,b)=1,
  \]
  heißen $a$ und $b$ \textbf{teilerfremd}.

  Für $a\in\bbN$ setzen wir
  \[
    \operatorname{ggT}(a,0)
    =\operatorname{ggT}(0,a):=a
  \]
  sowie
  \[
    \operatorname{ggT}(0,0):=0.
  \]

  Das \textbf{kleinste gemeinsame Vielfache}
  \[
    \operatorname{kgV}(a,b)
  \]
  ist die kleinste natürliche Zahl,
  die sowohl durch $a$ als auch durch $b$
  teilbar ist.
}


\mlemma{Eigenschaften von ggT und kgV}{
  Für alle $a,b\in\bbN$ gilt:
  \begin{enumerate}
    \item
      \[
        ab=\operatorname{ggT}(a,b)
        \cdot\operatorname{kgV}(a,b).
      \]

    \item
      \[
        \operatorname{ggT}(a,b)
        =\operatorname{ggT}(a\bmod b,b).
      \]
  \end{enumerate}
}

\begin{myproof}
  \textbf{1.}

  Setze
  \[
    d:=\operatorname{ggT}(a,b),
    \qquad a=du,\quad b=dv.
  \]
  Dann sind $u$ und $v$ teilerfremd.

  Die Zahl
  \[
    duv=\frac{ab}{d}
  \]
  ist ein gemeinsames Vielfaches von
  $a$ und $b$.

  Ist $m$ ein beliebiges gemeinsames
  Vielfaches, so existieren $s,t\in\bbN$
  mit
  \[
    m=as=bt.
  \]

  Daraus folgt
  \[
    us=vt.
  \]

  Wegen $\operatorname{ggT}(u,v)=1$
  gilt $v\mid s$. Somit ist
  \[
    duv\mid m.
  \]

  Daher ist $duv$ das kleinste gemeinsame
  Vielfache und es folgt
  \[
    \operatorname{kgV}(a,b)=\frac{ab}{d}.
  \]

  Also
  \[
    ab=\operatorname{ggT}(a,b)
       \operatorname{kgV}(a,b).
  \]

  \medskip
  \textbf{2.}

  Nach der Division mit Rest gilt
  \[
    a=b\left\lfloor\frac ab\right\rfloor
      +(a\bmod b).
  \]

  Ist $d$ ein gemeinsamer Teiler von
  $a$ und $b$, so teilt $d$ auch
  \[
    a-b\left\lfloor\frac ab\right\rfloor
    =a\bmod b.
  \]

  Umgekehrt teilt jeder gemeinsame Teiler
  von $a\bmod b$ und $b$ auch
  \[
    (a\bmod b)
    +b\left\lfloor\frac ab\right\rfloor
    =a.
  \]

  Somit besitzen $(a,b)$ und
  $(a\bmod b,b)$ dieselben gemeinsamen
  Teiler und daher denselben ggT.
\end{myproof}


\dfn{Euklidischer Algorithmus}{
  \textbf{Eingabe:} $a,b\in\bbN$.

  \textbf{Ausgabe:} $\operatorname{ggT}(a,b)$.

  \begin{enumerate}
    \item Solange $a>0$ und $b>0$:
      \begin{itemize}
        \item Falls $a<b$, setze
          \[
            b\leftarrow b\bmod a.
          \]
        \item Andernfalls setze
          \[
            a\leftarrow a\bmod b.
          \]
      \end{itemize}

    \item Gib $\max\{a,b\}$ aus.
  \end{enumerate}
}

\nt{
  Bei jedem Schritt bleibt der größte
  gemeinsame Teiler unverändert.

  Gleichzeitig wird mindestens eine der
  beiden Zahlen kleiner. Da alle Werte
  nichtnegative ganze Zahlen sind,
  terminiert der Algorithmus.

  Sobald einer der Werte $0$ ist,
  entspricht der andere dem ggT.
}


\ex{Euklidischer Algorithmus}{
  Wir berechnen
  \[
    \operatorname{ggT}(6314,2800).
  \]

  Durch wiederholte Division mit Rest:
  \begin{align*}
    6314&=2\cdot2800+714,\\
    2800&=3\cdot714+658,\\
    714&=1\cdot658+56,\\
    658&=11\cdot56+42,\\
    56&=1\cdot42+14,\\
    42&=3\cdot14+0.
  \end{align*}

  Somit folgt
  \[
    \boxed{\operatorname{ggT}(6314,2800)=14.}
  \]
}


\subsection{Fibonacci-Zahlen}

Für eine genauere Laufzeitanalyse des
euklidischen Algorithmus betrachten wir
die \textbf{Fibonacci-Zahlen}.

\dfn{Fibonacci-Zahlen}{
  Die Fibonacci-Folge $(F_n)_{n\in\bbN_0}$
  wird rekursiv definiert durch
  \[
    F_0:=0,\qquad F_1:=1
  \]
  und
  \[
    F_{n+1}:=F_n+F_{n-1}
    \qquad(n\ge1).
  \]

  Die ersten Folgenglieder lauten
  \[
    0,1,1,2,3,5,8,13,21,34,55,\dots
  \]
}


\mlemma{Explizite Formel der Fibonacci-Zahlen}{
  Für alle $n\in\bbN_0$ gilt
  \[
    \boxed{
      F_n=
      \frac1{\sqrt5}
      \left(
        \left(\frac{1+\sqrt5}{2}\right)^n
        -
        \left(\frac{1-\sqrt5}{2}\right)^n
      \right).
    }
  \]
}

\begin{myproof}
  Wir führen vollständige Induktion über $n$.

  \textbf{Induktionsanfang:}

  Für $n=0$ gilt
  \[
    F_0=\frac{1-1}{\sqrt5}=0.
  \]

  Für $n=1$ erhalten wir
  \[
    F_1
    =\frac1{\sqrt5}
     \left(
       \frac{1+\sqrt5}{2}
       -\frac{1-\sqrt5}{2}
     \right)
    =1.
  \]

  \textbf{Induktionsschritt:}

  Setze
  \[
    \varphi:=\frac{1+\sqrt5}{2},
    \qquad
    \psi:=\frac{1-\sqrt5}{2}.
  \]

  Es gilt
  \[
    \varphi^2=\varphi+1,
    \qquad
    \psi^2=\psi+1.
  \]

  Angenommen, die Formel gilt für
  $n-1$ und $n-2$.

  Dann folgt für $n\ge2$:
  \begin{align*}
    F_n
    &=F_{n-1}+F_{n-2}\\
    &=\frac1{\sqrt5}
      \left(
        \varphi^{n-1}-\psi^{n-1}
        +\varphi^{n-2}-\psi^{n-2}
      \right)\\
    &=\frac1{\sqrt5}
      \left(
        \varphi^{n-2}(\varphi+1)
        -\psi^{n-2}(\psi+1)
      \right)\\
    &=\frac1{\sqrt5}
      \left(\varphi^n-\psi^n\right).
  \end{align*}

  Damit gilt die Formel für alle
  $n\in\bbN_0$.
\end{myproof}


\subsection{Laufzeit des euklidischen Algorithmus}

\mlemma{Fibonacci-Abschätzung}{
  Seien $a>b>0$ und durchlaufe der
  euklidische Algorithmus $k\ge1$
  Iterationen.

  Dann gilt
  \[
    \boxed{
      a\ge F_{k+2},
      \qquad
      b\ge F_{k+1}.
    }
  \]
}

\begin{myproof}
  Wir führen Induktion über $k$.

  \textbf{Induktionsanfang:}

  Für $k=1$ gilt wegen $a>b>0$:
  \[
    a\ge2=F_3,
    \qquad
    b\ge1=F_2.
  \]

  \textbf{Induktionsschritt:}

  Sei $k\ge2$ und die Aussage für
  $k-1$ Iterationen bereits bewiesen.

  Nach der ersten Iteration betrachten
  wir die Zahlen
  \[
    b,\qquad a\bmod b.
  \]

  Da noch $k-1$ Iterationen verbleiben,
  gilt nach Induktionsvoraussetzung
  \[
    b\ge F_{k+1},
    \qquad
    a\bmod b\ge F_k.
  \]

  Wegen $a>b$ gilt
  \[
    \left\lfloor\frac ab\right\rfloor\ge1.
  \]

  Somit folgt
  \begin{align*}
    a
    &=\left\lfloor\frac ab\right\rfloor b
      +(a\bmod b)\\
    &\ge b+(a\bmod b)\\
    &\ge F_{k+1}+F_k\\
    &=F_{k+2}.
  \end{align*}

  Damit ist die Behauptung bewiesen.
\end{myproof}


\thm{Satz von Lamé}{
  Seien $a\ge b$ und
  \[
    b<F_{k+1}
  \]
  für ein $k\in\bbN$.

  Dann benötigt der euklidische Algorithmus
  \textbf{weniger als $k$ Iterationen}.
}

\begin{myproof}
  Falls $b=0$, benötigt der Algorithmus
  keine Iteration.

  Falls $a=b>0$, benötigt er genau eine
  Iteration. Für $k\ge2$ ist dies weniger
  als $k$; für $k=1$ ist die Voraussetzung
  $b<F_2=1$ nicht erfüllt.

  Sei daher $a>b>0$.

  Angenommen, der Algorithmus benötige
  mindestens $k$ Iterationen.

  Nach dem vorherigen Lemma würde dann
  gelten
  \[
    b\ge F_{k+1}.
  \]

  Dies widerspricht der Voraussetzung
  \[
    b<F_{k+1}.
  \]

  Also benötigt der Algorithmus weniger
  als $k$ Iterationen.
\end{myproof}


\nt{
  Der Satz von Lamé ist \textbf{bestmöglich}.

  Für zwei aufeinanderfolgende Fibonacci-Zahlen
  \[
    a=F_{k+2},
    \qquad
    b=F_{k+1}
  \]
  benötigt der euklidische Algorithmus
  genau $k$ Iterationen.

  Denn wegen
  \[
    F_{k+2}=F_{k+1}+F_k
  \]
  werden die Eingaben bei jeder Iteration
  auf das vorherige Paar von
  Fibonacci-Zahlen zurückgeführt.
}


\cor{Laufzeit des euklidischen Algorithmus}{
  Seien $a\ge b>1$ und
  \[
    \varphi:=\frac{1+\sqrt5}{2}.
  \]

  Dann benötigt der euklidische Algorithmus
  höchstens
  \[
    \boxed{\lceil\log_\varphi b\rceil}
  \]
  Iterationen.

  Die Gesamtlaufzeit beträgt
  \[
    \boxed{O((\log a)^3).}
  \]
}

\begin{myproof}
  Zunächst zeigen wir durch Induktion,
  dass für alle $n\ge3$ gilt:
  \[
    F_n>\varphi^{n-2}.
  \]

  \textbf{Induktionsanfang:}

  Für $n=3$ gilt
  \[
    F_3=2>\varphi.
  \]

  \textbf{Induktionsschritt:}

  Für $n\ge4$ folgt aus den
  Induktionsvoraussetzungen
  \begin{align*}
    F_{n+1}
    &=F_n+F_{n-1}\\
    &>\varphi^{n-2}
      +\varphi^{n-3}\\
    &=\varphi^{n-3}(\varphi+1)\\
    &=\varphi^{n-1},
  \end{align*}
  da $\varphi^2=\varphi+1$.

  Somit gilt die Behauptung.

  \medskip
  Setze nun
  \[
    k:=1+\lceil\log_\varphi b\rceil.
  \]

  Dann ist $k\ge2$ und
  \[
    b\le\varphi^{k-1}<F_{k+1}.
  \]

  Nach dem Satz von Lamé benötigt
  der Algorithmus weniger als $k$,
  also höchstens
  \[
    \lceil\log_\varphi b\rceil
  \]
  Iterationen.

  \medskip
  \textbf{Gesamtlaufzeit:}

  In jeder Iteration wird eine Division
  mit Rest berechnet.

  Für Zahlen mit höchstens
  $O(\log a)$ Binärstellen benötigt
  die schriftliche Division höchstens
  \[
    O((\log a)^2)
  \]
  Rechenschritte.

  Da es höchstens $O(\log a)$
  Iterationen gibt, folgt insgesamt
  \[
    O(\log a)\cdot O((\log a)^2)
    =\boxed{O((\log a)^3)}.
  \]
\end{myproof}


\nt{
  Die Zahl
  \[
    \varphi=\frac{1+\sqrt5}{2}
    \approx1{,}618
  \]
  heißt \textbf{goldener Schnitt}.

  Da
  \[
    \left|\frac{1-\sqrt5}{2}\right|<1,
  \]
  folgt aus der expliziten Fibonacci-Formel
  die Näherung
  \[
    F_n\approx\frac{\varphi^n}{\sqrt5}.
  \]

  Die Fibonacci-Zahlen wachsen also
  exponentiell. Deshalb ist die Anzahl
  der Iterationen des euklidischen
  Algorithmus nur logarithmisch
  in der Größe der Eingabe.
}
